\par\noindent\begin{minipage}{\linewidth}
\small
\setlength{\tabcolsep}{4pt}
\renewcommand{\arraystretch}{1.13}
\captionof{table}{Oposición persistente de los mismos modelos con medidas y conjuntos alternativos.}\label{tab:S15}
\begin{tabular}{@{}>{\raggedright\arraybackslash}p{3.55cm}>{\raggedright\arraybackslash}p{1.7cm}>{\raggedright\arraybackslash}p{1.8cm}>{\raggedright\arraybackslash}p{3.15cm}>{\raggedright\arraybackslash}p{1.5cm}>{\raggedright\arraybackslash}p{3.05cm}@{}}
\toprule
\textbf{Conjunto de modelos} & \textbf{Modelos} & \textbf{Apertura} & \textbf{Apertura + alternativas} & \textbf{PR} & \textbf{PR + alternativas} \\
\midrule
Los diez & 10 & 262 & 225 & 221 & 196 \\
Seis de similitud & 6 & 231 & 190 & 160 & 121 \\
Sin SPECTER & 9 & 234 & 182 & 212 & 181 \\
Sin SPECTER2 & 9 & 251 & 205 & 211 & 182 \\
Sin SciNCL & 9 & 259 & 217 & 215 & 190 \\
Sin SciBERT & 9 & 258 & 222 & 218 & 194 \\
Sin BERT & 9 & 261 & 223 & 211 & 183 \\
Sin MPNet & 9 & 262 & 223 & 218 & 192 \\
Sin MiniLM & 9 & 259 & 223 & 221 & 196 \\
Sin PubMedBERT & 9 & 261 & 220 & 206 & 188 \\
Sin BioBERT & 9 & 259 & 217 & 212 & 181 \\
Sin SimCSE & 9 & 259 & 217 & 217 & 191 \\
\bottomrule
\end{tabular}
\end{minipage}\par
\par\smallskip{\small Todos los recuentos son sobre 325 y con magnitud mínima cero. Las alternativas deben conservar las direcciones de los mismos modelos opuestos; no se sustituyen por otros entre medidas. Las alternativas angulares tienen 26 condiciones guardadas y las espectrales tres. Los seis de similitud son SPECTER, SPECTER2, SciNCL, MPNet, MiniLM y SimCSE. Menos modelos ofrecen menos oportunidades de oposición; este conjunto no aísla un efecto del entrenamiento. Los CSV incluyen los cuatro cortes de magnitud y la identidad de los modelos que sustentan la oposición.\par}
